% sections/results.tex -- Section 7 (paper_v2, cut in C5 from the full version's). The details that
% left it (the paths' ranks, the delta check, the user groups, more paths, the path weights and
% preferences, HAN against metapath2vec) are in the appendix sections on the search and the full
% results, which become the supplementary document in B5. B3 (25 November) brought in the reviewed
% Vietnamese version (paper_v2_vn, L10): the random draws as the control (7.2, with the cheapest
% choice), cost against effectiveness with the search's GPU hours as the yardstick (7.3), why the
% searched sets lose (end of 7.4) and why GAT leads (7.5), both from the Discussion, and a few
% numbers from the appendix in the text. Every number is \val{} from numbers.tex. Tables: T4
% (tab:main), T7 (tab:cost); figures: F3 (fig:sets), F4 (fig:weights).
\section{Results}
\label{sec:results}

All numbers are test results unless stated otherwise.

\subsection{Searched against manual sets}
\label{sec:res-main}

On the metric fixed in advance, \code{a1} loses to the manual set of its size in both models
(Table~\ref{tab:main}). In HAN it trails M5 by \val{han-4-hit-diffabs} (\val{han-4-hit-relabs}\%,
\val{han-4-hit-pt}) and M6 by \val{han-6-hit-diffabs} (\val{han-6-hit-relabs}\%,
\val{han-6-hit-pt}), \val{han-4-hit-sdx} and \val{han-6-hit-sdx} times the larger seed standard
deviation. In metapath2vec the gaps are wider: \val{mtv-4-hit-diffabs} (\val{mtv-4-hit-relabs}\%)
and \val{mtv-6-hit-diffabs} (\val{mtv-6-hit-relabs}\%). On full ranking HAN's sets tie
(\val{han-4-fhit-pt} and \val{han-6-fhit-pt}). In metapath2vec, \code{a1} ties at four paths
(\val{mtv-4-fhit-pt}) and loses at six (\val{mtv-6-fhit-pt}). In HAN the loss holds for \code{a2}
and \code{c} too (Figure~\ref{fig:sets}) and in every user group (differences of
\val{han-6-hit-d3-diff} to \val{han-4-hit-d1-diff}).

% t4_main.tex -- T4 of the plan: the main comparison (results.csv via numbers.tex; Welch tests
% computed by make_numbers.py).
\begin{widetable}[htbp]
  \centering
  \small
  \setlength{\tabcolsep}{4pt}
  \caption{Searched against manual sets: test results, mean $\pm$ standard deviation over 3
    training seeds. $\Delta$: the searched set minus the manual set, with Welch's two-sided $p$.
    Sampled: each test edge ranked against \val{split-negatives} negatives, with Hit@10 the metric
    fixed in advance. Full: against
    every product the user has no edge to, with Recall@20 averaged over users.}
  \label{tab:main}
  \begin{tabular*}{\linewidth}{@{\extracolsep{\fill}}llllll@{}}
    \toprule
    & & \multicolumn{2}{c}{Sampled} & \multicolumn{2}{c}{Full ranking} \\
    \cmidrule(lr){3-4}\cmidrule(l){5-6}
    Model & Set & Hit@10 & NDCG@10 & Hit@10 & Recall@20 \\
    \midrule
    HAN & \code{a1}, 4 paths & \val{han-f4a1-hit-mean} $\pm$ \val{han-f4a1-hit-sd}
      & \val{han-f4a1-ndcg-mean} & \val{han-f4a1-fhit-mean} $\pm$ \val{han-f4a1-fhit-sd}
      & \val{han-f4a1-frec-mean} \\
    & M5 & \val{han-m5-hit-mean} $\pm$ \val{han-m5-hit-sd} & \val{han-m5-ndcg-mean}
      & \val{han-m5-fhit-mean} $\pm$ \val{han-m5-fhit-sd} & \val{han-m5-frec-mean} \\
    & $\Delta$ ($p$) & \val{han-4-hit-diff} (\val{han-4-hit-p})
      & \val{han-4-ndcg-diff} (\val{han-4-ndcg-p}) & \val{han-4-fhit-diff} (\val{han-4-fhit-p})
      & \val{han-4-frec-diff} (\val{han-4-frec-p}) \\
    \addlinespace
    & \code{a1}, 6 paths & \val{han-f6a1-hit-mean} $\pm$ \val{han-f6a1-hit-sd}
      & \val{han-f6a1-ndcg-mean} & \val{han-f6a1-fhit-mean} $\pm$ \val{han-f6a1-fhit-sd}
      & \val{han-f6a1-frec-mean} \\
    & M6 & \val{han-m6-hit-mean} $\pm$ \val{han-m6-hit-sd} & \val{han-m6-ndcg-mean}
      & \val{han-m6-fhit-mean} $\pm$ \val{han-m6-fhit-sd} & \val{han-m6-frec-mean} \\
    & $\Delta$ ($p$) & \val{han-6-hit-diff} (\val{han-6-hit-p})
      & \val{han-6-ndcg-diff} (\val{han-6-ndcg-p}) & \val{han-6-fhit-diff} (\val{han-6-fhit-p})
      & \val{han-6-frec-diff} (\val{han-6-frec-p}) \\
    \midrule
    metapath2vec & \code{a1}, 4 paths & \val{mtv-f4a1-hit-mean} $\pm$ \val{mtv-f4a1-hit-sd}
      & \val{mtv-f4a1-ndcg-mean} & \val{mtv-f4a1-fhit-mean} $\pm$ \val{mtv-f4a1-fhit-sd}
      & \val{mtv-f4a1-frec-mean} \\
    & M5 & \val{mtv-m5-hit-mean} $\pm$ \val{mtv-m5-hit-sd} & \val{mtv-m5-ndcg-mean}
      & \val{mtv-m5-fhit-mean} $\pm$ \val{mtv-m5-fhit-sd} & \val{mtv-m5-frec-mean} \\
    & $\Delta$ ($p$) & \val{mtv-4-hit-diff} (\val{mtv-4-hit-p})
      & \val{mtv-4-ndcg-diff} (\val{mtv-4-ndcg-p}) & \val{mtv-4-fhit-diff} (\val{mtv-4-fhit-p})
      & \val{mtv-4-frec-diff} (\val{mtv-4-frec-p}) \\
    \addlinespace
    & \code{a1}, 6 paths & \val{mtv-f6a1-hit-mean} $\pm$ \val{mtv-f6a1-hit-sd}
      & \val{mtv-f6a1-ndcg-mean} & \val{mtv-f6a1-fhit-mean} $\pm$ \val{mtv-f6a1-fhit-sd}
      & \val{mtv-f6a1-frec-mean} \\
    & M6 & \val{mtv-m6-hit-mean} $\pm$ \val{mtv-m6-hit-sd} & \val{mtv-m6-ndcg-mean}
      & \val{mtv-m6-fhit-mean} $\pm$ \val{mtv-m6-fhit-sd} & \val{mtv-m6-frec-mean} \\
    & $\Delta$ ($p$) & \val{mtv-6-hit-diff} (\val{mtv-6-hit-p})
      & \val{mtv-6-ndcg-diff} (\val{mtv-6-ndcg-p}) & \val{mtv-6-fhit-diff} (\val{mtv-6-fhit-p})
      & \val{mtv-6-frec-diff} (\val{mtv-6-frec-p}) \\
    \midrule
    Popularity & -- & \val{pop-hit} & \val{pop-ndcg} & \val{pop-fhit} & \val{pop-frec} \\
    \bottomrule
  \end{tabular*}
\end{widetable}


\subsection{The random control}
\label{sec:res-random}

To judge the gap between search and manual design against how much the choice of paths matters at
all, Figure~\ref{fig:sets} places the searched and the manual sets among all 24, with the random
draws as the control. Of the eight sets of each size drawn from the same \val{search-cands}
candidates (three searched and five random), HAN puts the searched sets \val{han-4-pos-hit} at four
paths and \val{han-6-pos-hit} at six, as validation BPR does. On full ranking they come
\val{han-4-pos-fhit} at four paths and \val{han-6-pos-fhit} at six. The margins lie within seed
noise: the best searched set reaches \val{han-f4c-hit-s1} at four paths and \val{han-f6a1-hit-s1}
at six, against \val{han-r4-hit-min} and \val{han-r6-hit-min} for the weakest draws. So the places
show that the search does not beat its own candidates, not that every draw beats it. The manual
sets do about as well as an average draw: M5 reaches \val{han-m5-hit-mean} against a mean of
\val{han-r4-hit-mean} for the four-path draws (range \val{han-r4-hit-min}--\val{han-r4-hit-max}),
and M6 \val{han-m6-hit-mean} against \val{han-r6-hit-mean}
(\val{han-r6-hit-min}--\val{han-r6-hit-max}).

% f3_sets.tex -- F3 of the plan: HAN test Hit@10 of all 24 sets by family. Data: f3_*.dat,
% written by make_numbers.py from results.csv (seed 1 of every set).
\begin{widefigure}[htbp]
  \centering
  \begin{tikzpicture}
    \begin{axis}[
        width=0.9\linewidth, height=5.4cm,
        xmin=0.45, xmax=3.6, ymin=0.612, ymax=0.712,
        xtick={1,2,3}, xticklabels={2--3 paths, 4 paths, 6 paths},
        ylabel={HAN test Hit@10 (sampled)},
        yticklabel style={/pgf/number format/fixed, /pgf/number format/precision=2},
        ytick={0.62,0.64,0.66,0.68,0.70},
        tick label style={font=\scriptsize}, label style={font=\footnotesize},
        ymajorgrids, grid style={black!8}, axis line style={black!60}, tick style={black!60},
        legend style={font=\scriptsize, at={(0.02,0.98)}, anchor=north west, draw=none,
          fill=white, fill opacity=0.85, text opacity=1},
        legend cell align=left]
      \draw[black!20] (axis cs:1.5,0.612) -- (axis cs:1.5,0.712);
      \draw[black!20] (axis cs:2.5,0.612) -- (axis cs:2.5,0.712);
      \addplot[only marks, mark=*, mark size=2.6pt, mark options={fill=cSearch, draw=white,
          line width=0.4pt}] table[x=x, y=hit] {figures/f3_searched.dat};
      \addlegendentry{searched (LMSPS)}
      \addplot[only marks, mark=square*, mark size=2.4pt, mark options={fill=cManual, draw=white,
          line width=0.4pt}] table[x=x, y=hit] {figures/f3_manual.dat};
      \addlegendentry{manual}
      \addplot[only marks, mark=triangle*, mark size=3pt, mark options={fill=cRandom, draw=white,
          line width=0.4pt}] table[x=x, y=hit] {figures/f3_random.dat};
      \addlegendentry{random draws}
      \addplot[only marks, mark=none, point meta=explicit symbolic,
          visualization depends on={value \thisrow{anchor}\as\labelanchor},
          nodes near coords, nodes near coords style={font=\scriptsize, text=black!80,
            anchor=\labelanchor, inner sep=2.5pt}]
        table[x=x, y=y, meta=label] {figures/f3_labels.dat};
    \end{axis}
  \end{tikzpicture}
  \Description{A dot plot of HAN's test Hit@10 for all 24 path sets grouped by set size, with one marker shape and colour per family. The searched sets sit at the low end of the sets drawn from the candidates, and the manual sets sit among the random draws.}
  \caption{HAN's test Hit@10 for all 24 path sets, one seed each (the main sets' standard
    deviation over three seeds is at most \val{han-m6-hit-sd}). At this seed the searched sets come last among the sets
    drawn from the same \val{search-cands} candidates at both sizes, though by margins within seed noise. M5 and M6 sit near
    the middle of the random draws. For scale: popularity \val{pop-hit}, GAT \val{gat-hit-mean}.}
  \label{fig:sets}
\end{widefigure}


The choice of paths thus moves HAN little: all \val{han-4-nsets} four-path sets lie within
\val{han-4-span} Hit@10 of each other, and all \val{han-6-nsets} six-path sets within
\val{han-6-span}. Choosing by validation works: over the 24 sets, validation BPR and test Hit@10
have a Spearman correlation of \val{han-valtest-rho} (\val{han-valfull-rho} with full-ranking
Hit@10). The validation pick is a random draw and reaches \val{han-r6d2-hit-s1}, which no other set
beats on the metric fixed in advance. The cheapest choice is thus to draw a few sets and keep the
one with the best validation score. In metapath2vec the comparison with chance is mixed: the
searched sets come \val{mtv-4-pos-hit} of eight at four paths (above the average draw) but
\val{mtv-6-pos-hit} at six (below it). The two models also prefer different sets: over the 23 sets
with a path through users, their Hit@10 values have a Spearman correlation of only
\val{han-mtv-rho}.

\subsection{Cost against effectiveness}
\label{sec:res-cost}

The hand design used no GPU, and its time was not recorded either (Table~\ref{tab:cost}). Earlier
work also describes this cost only in words \citep[][App.~D]{han2020genetic,li2024longrange}. LMSPS
even leaves the time for hand design out of its efficiency comparison, together with its own search
time \citep[App.~E.1]{li2024longrange}. The search's cost still gives us a yardstick: the
\val{search-runs-total} search runs took about \val{search-total-h} GPU hours, not counting the time
to adapt the code. Hand design is cheaper than the search as long as it takes less than this. We
think this is plausible on a schema with only three node types, where the manual sets were designed
from four principles in one session and the density checks took a few CPU minutes. Meanwhile, the
search gives sets that trail the manual ones by \val{han-4-hit-relabs}\% and
\val{han-6-hit-relabs}\% on the metric fixed in advance and only tie them on full ranking. On this
schema the search's cost has therefore not bought HAN better results than hand design. The random
control costs nothing and also does at least as well as both. On full ranking the random draws
average at least as high, though with a wider spread (\val{han-r4-fhit-min}--\val{han-r4-fhit-max}
and \val{han-r6-fhit-min}--\val{han-r6-fhit-max}).

% t7_cost.tex -- T7 of the plan: cost and performance of the three ways to choose HAN's paths.
% Search cost from the 60 search logs (150 epochs at the logged epoch time plus the feature
% propagation); results from results.csv; all via numbers.tex. The cost of the hand design is
% stated in words (plan, D10): none of the reference papers measures it.
\begin{widetable}[htbp]
  \centering
  \small
  \caption{What each way of choosing HAN's meta-paths costs and what it gives: test Hit@10,
    sampled and full ranking. By hand and by search: mean over 3 seeds. At random: mean of 5 draws,
    one seed each. The relative gap is to the manual set of the same size. Training HAN on the
    chosen set costs the same for all three, a median \val{han-run-min} GPU minutes per run.}
  \label{tab:cost}
  \begin{tabularx}{\linewidth}{@{}lLlll@{}}
    \toprule
    & & \multicolumn{2}{l}{Sampled Hit@10} & Full Hit@10 \\
    \cmidrule(lr){3-4}\cmidrule(l){5-5}
    Chosen & Cost of choosing & 4 / 6 paths & vs.\ by hand & 4 / 6 paths \\
    \midrule
    By hand (M5, M6) & no GPU: one AI-assisted design session (time not recorded) and density
      checks on a small graph in CPU minutes & \val{han-m5-hit-mean} / \val{han-m6-hit-mean} & --
      & \val{han-m5-fhit-mean} / \val{han-m6-fhit-mean} \\
    By search (\code{a1}) & \val{search-runs} runs per size, a median \val{search-run-min-4} /
      \val{search-run-min-6} GPU minutes each: \val{search-total-h} GPU hours
      & \val{han-f4a1-hit-mean} / \val{han-f6a1-hit-mean}
      & \val{han-4-hit-rel}\% / \val{han-6-hit-rel}\%
      & \val{han-f4a1-fhit-mean} / \val{han-f6a1-fhit-mean} \\
    At random (5 draws) & none & \val{han-r4-hit-mean} / \val{han-r6-hit-mean}
      & \val{han-r4-vs-manual-rel}\% / \val{han-r6-vs-manual-rel}\%
      & \val{han-r4-fhit-mean} / \val{han-r6-fhit-mean} \\
    \bottomrule
  \end{tabularx}
\end{widetable}


\subsection{What HAN does with the paths}
\label{sec:res-weights}

Figure~\ref{fig:weights} compares the search's preference for each path with HAN's weight for it.
The two agree on paths through users. \code{PUPUP} (which links products through two users) is in
all \val{searched-with-pupup} searched sets and \code{PUPPUP} in \val{searched-with-puppup}. Both
rank near the top of the search's candidates. HAN never drops them or \code{PUP} and weights all
three above an equal share. They disagree on categories. Each searched set holds
\val{searched-empty-min} to \val{searched-empty-max} near-empty category paths. The search ranks
these high, but HAN keeps only one of them per set with a weight of \val{han-empty-kept-min} to
\val{han-empty-kept-max} (\val{han-empty-kept-f4c} for \code{c} at four paths), so the six-path
searched sets work like four-path sets. The dense category paths of the manual sets rank low in the
search but get about an equal share or more in HAN (median \val{kind-cat-weight}). No searched set
holds a category path that reaches most products.

% f4_weights.tex -- F4 of the plan: HAN's path weights against the search's preference, by kind
% of path. Data: f4_*.dat, written by make_numbers.py (search records; HAN seed-1 logs).
\begin{figure}[htbp]
  \centering
  \begin{tikzpicture}
    \begin{axis}[
        width=\linewidth, height=5.6cm,
        xmin=0, xmax=33, ymin=-0.1, ymax=2.05, xtick={1,8,16,24,32},
        xlabel={Search's preference: mean $\alpha$ rank (1 = best of \val{search-cands})},
        ylabel={HAN's weight (1 = an equal share)},
        tick label style={font=\scriptsize}, label style={font=\footnotesize},
        ymajorgrids, grid style={black!8}, axis line style={black!60}, tick style={black!60},
        legend style={font=\scriptsize, at={(0.5,1.02)}, anchor=south, draw=none,
          legend columns=2, /tikz/every even column/.append style={column sep=6pt}},
        legend cell align=left]
      \draw[black!45, dashed] (axis cs:0,1) -- (axis cs:33,1);
      \addplot[only marks, mark=*, mark size=2.6pt, mark options={fill=cSearch, draw=white,
          line width=0.4pt}] table[x=rank, y=weight] {figures/f4_users.dat};
      \addlegendentry{through users (\val{kind-users-paths} paths)}
      \addplot[only marks, mark=square*, mark size=2.4pt, mark options={fill=cManual, draw=white,
          line width=0.4pt}] table[x=rank, y=weight] {figures/f4_empty.dat};
      \addlegendentry{near-empty category (\val{kind-empty-paths})}
      \addplot[only marks, mark=triangle*, mark size=3pt, mark options={fill=cRandom, draw=white,
          line width=0.4pt}] table[x=rank, y=weight] {figures/f4_cat.dat};
      \addlegendentry{other category (\val{kind-cat-paths})}
      \addplot[only marks, mark=diamond*, mark size=3pt, mark options={fill=cOther, draw=white,
          line width=0.4pt}] table[x=rank, y=weight] {figures/f4_pp.dat};
      \addlegendentry{products only (\val{kind-pp-paths})}
      \addplot[only marks, mark=none, point meta=explicit symbolic,
          visualization depends on={value \thisrow{anchor}\as\labelanchor},
          nodes near coords, nodes near coords style={font=\scriptsize\ttfamily, text=black!80,
            anchor=\labelanchor, inner sep=2.5pt}]
        table[x=x, y=y, meta=label] {figures/f4_labels.dat};
    \end{axis}
  \end{tikzpicture}
  \Description{A scatter plot with one point per path: the search's mean preference for the path against HAN's mean weight for it. Paths through users are preferred by both. Near-empty category paths are preferred by the search but get almost no weight in HAN. Dense category paths are ranked low by the search but used by HAN.}
  \caption{What HAN does with the paths: each point is a path, at the search's mean preference for
    it (both set sizes) and HAN's mean weight for it over the sets that hold it (seed 1). The two
    agree on paths through users. They disagree on the near-empty category paths, which the search
    ranks high and HAN mostly drops (\val{kind-empty-dropped} of \val{kind-empty-uses} uses get a
    weight below \val{weight-drop}). They also disagree on the dense category paths, which the search
    ranks low and HAN uses.}
  \label{fig:weights}
\end{figure}


Among the sets of four and of six paths, that lack separates the searched sets from all the others.
At both sizes, every set with a category path that reaches almost every product beats every set
without one: \val{han-4-dense-min}--\val{han-4-dense-max} against
\val{han-4-nodense-min}--\val{han-4-nodense-max} at four paths (the closest pair within seed noise)
and \val{han-6-dense-min}--\val{han-6-dense-max} against \val{han-6-nodense-max} at six. Across the
\val{pref-cand-n} sets drawn from the candidates, the more the search likes a set's paths, the
worse HAN does with it: test Hit@10 and the paths' mean $\alpha$ rank (lower for a stronger
preference) have a Spearman correlation of \val{pref-cand-rho} (\val{pref-cand-pt}). This relation
rests on the searched sets: within the ten random draws alone, the coefficient is only
\val{pref-random-rho} (\val{pref-random-pt}).

\paragraph{Where the loss sits.} A single step along the category hierarchy (\code{PCCP}) links
\val{bg-pccp-reach-products-pct}\% of the products to another product. These are mostly DVDs and
videos, which hold \val{bg-pccp-reach-test-pct}\% of the test edges. On the products \code{PCCP}
does not reach, the searched sets trail the sets with a dense category path by \val{bg-4-gapabs-no}
Hit@10 at four paths and by \val{bg-6-gapabs-no} at six (seed~1). On the products it reaches, they
trail by only \val{bg-4-gapabs-reach} and \val{bg-6-gapabs-reach}. The loss thus sits mostly on the
books and music for which the searched sets' category paths are empty.

\paragraph{Why the searched sets lose.} The check above points to the near-empty category paths. In
the search, a product that such a path cannot reach gets an all-zero feature for it. The zero still
carries information: the products that one hierarchy step does reach (mostly DVDs and videos) hold
\val{pccp-reach-train-pct}\% of the training edges but only \val{bg-pccp-reach-products-pct}\% of
the products and hence of the random negatives. The empty paths thus reveal the product group and
part of a product's popularity. In HAN this cue is mostly gone: every product keeps its own features
through a self-loop in every graph, so an empty channel only repeats them. HAN drops most of these
channels. The searched sets then lack a category path that reaches books and music, and that is
where they lose most. The search may also make less of a dense path than HAN, as it averages the
path's features over all endpoints, while HAN attends over \val{han-set-topk} selected neighbours.
However, the product-group check only locates the loss and does not show that a dense category path
would remove it.

\subsection{Other models}
\label{sec:res-models}

GAT beats both meta-path models on every metric and in every user group. It reaches a Hit@10 of
\val{gat-hit-mean}, against \val{han-r6d2-hit-s1} for HAN's validation pick and
\val{mtv-m5-hit-mean} for metapath2vec's. On full ranking it reaches \val{gat-fhit-mean}, against
\val{han-r6d2-fhit-s1} for HAN. Its lead over HAN is largest for users with few training edges:
\val{gat-lead-d1} for 1--5 edges, \val{gat-lead-d2} for 6--20 and \val{gat-lead-d3} for 21 or
more. Unlike the two meta-path models, GAT builds each user from the products the user reviewed,
while HAN and metapath2vec learn a free vector per user that gets little signal from a handful of
edges. This fits HGB's finding that a homogeneous GAT with proper inputs is a strong baseline
\citep{lv2021hgb}. It also suggests that here the way users are represented matters more than which
product-side meta-paths are chosen. Our design leaves out user-side meta-paths, which would test
this.

HAN beats metapath2vec on sampled Hit@10 in all \val{han-wins-sampled} sets, but metapath2vec wins on
full ranking in \val{mtv-wins-full}, so the order of the two depends on the metric. Every model
clears the popularity floor (\val{pop-hit} sampled, \val{pop-fhit} on full ranking) except
metapath2vec on M4, which reaches \val{mtv-m4-hit-s1}: that set has no path through users, so the
users stay untrained.

\subsection{Sensitivity}
\label{sec:res-sensitivity}

With a path-weight warm-up of \val{wu5-warmup} epochs instead of \val{han-set-beta-warmup}
(seed~1, all 24 sets), HAN trains worse on average (\val{wu5-diff-mean} Hit@10, from
\val{wu5-diff-min} to \val{wu5-diff-max}), drops more paths (\val{wu5-dropped} weights below
\val{weight-drop}, against \val{han-dropped}) and orders the sets differently (Spearman
\val{wu5-rho} with the main runs). The searched sets then come \val{wu5-4-pos} of eight at both
sizes instead of last. At six paths the seed-1 comparison reverses: \code{a1} reaches
\val{wu5-f6a1-hit} against \val{wu5-m6-hit} for M6, which keeps only \val{wu5-m6-kept} of its six
paths. This check uses one seed and was not planned in advance, so it does not change the main
result, but it shows that the comparison depends on how HAN learns its path weights: our
conclusions hold for the \val{han-set-beta-warmup}-epoch warm-up.
